\begin{statementbox}
	\textbf{\textcolor{CStatement}{条件}}
	\begin{description}[style=nextline, leftmargin=2.8em, labelsep=0.5em, itemsep=0.35em]
		\item[\textbf{\textcolor{CStatement}{条件 1（公差非零）：}}]
			等差数列$\{a_n\}$的公差$d \neq 0$．
	\end{description}
	\textbf{\textcolor{CStatement}{结论}}
	\begin{description}[style=nextline, leftmargin=2.8em, labelsep=0.5em, itemsep=0.45em]
		\item[\textbf{\textcolor{CStatement}{结论一（最大值条件）：}}]
			\textbf{\textcolor{CStatement}{直观描述：}}
			当公差为负且项由正转负时，和取最大。
			\[
				d<0, a_n\ge 0, a_{n+1}\le 0
			\]
			此时$S_n$取得最大值。
		\item[\textbf{\textcolor{CStatement}{结论二（最小值条件）：}}]
			\textbf{\textcolor{CStatement}{直观描述：}}
			当公差为正且项由负转正时，和取最小。
			\[
				d>0, a_n\le 0, a_{n+1}\ge 0
			\]
			此时$S_n$取得最小值。
	\end{description}

	\medskip
	\textbf{\textcolor{CStatement}{推广结论（二次函数视角）：}}
	\textbf{\textcolor{CStatement}{直观描述：}}
	将前$n$项和写作关于$n$的二次函数，由开口方向确定最值类型。
	\[
		S_n = \frac{d}{2}n^2 + \Bigl(a_1-\frac{d}{2}\Bigr)n,
	\]
	开口由$\frac{d}{2}$符号决定：$d>0$开口向上（有最小值），$d<0$开口向下（有最大值）。最值点在对称轴$n=-\frac{a_1-\frac{d}{2}}{d}$附近取整。
\end{statementbox}
